\question[5]
Completa la tabella per i tre esempi del laboratorio: offset, struttura del payload e risultato ottenuto.
\begin{solution}
\begin{center}\small
\begin{tabular}{@{}llll@{}}
\toprule
Esempio & Offset & Payload & Risultato\\
\midrule
override   & 12 & \lstinline|padding + p32(valore)|          & variabile controllata\\
corruption & 32 & \lstinline|padding + p32(&funzione)|        & funzione nascosta eseguita\\
shellcode  & 76 & \lstinline|padding + indirizzo + shellcode| & shell interattiva\\
\bottomrule
\end{tabular}
\end{center}
Strumenti trasversali: \lstinline|cyclic| per l'offset, \lstinline|p32|/\lstinline|p64| per l'impacchettamento, \lstinline|process()| per l'invio.
\end{solution}
